\scanpage{50}
\mspage{24}
Tout morphisme de générisation se factorise en un composé de morphismes de générisation d'écart $1$. \deleted{de façon unique} \inserted{Le nb. de} façons de factoriser ainsi, à isom. près, est égal à $\delta!$, où $\delta$ est l'écart du morphisme de générisation.

Considérons, pour un MD-graphe donné $\underline{G}$, de type \inserted{numérique} $(g,\deleted{\unsure{I}}\nu)$, \deleted{les} \inserted{inclusions de} schémas modulaires \inserted{propres sur $S$}
\diagram{\begin{tikzcd}[ampersand replacement=\&]
\text{(57)}\quad\widehat M_{[\underline G]}\arrow[r,hook]\&\widehat M_{g,\nu}
\end{tikzcd}}
et aussi l'inclusion
\diagram{\begin{tikzcd}[ampersand replacement=\&]
\text{(58)}\quad\underbrace{(\widehat M_G,\Gamma^!)}_{\deleted{\unsure{\widehat M}}\unsure{C}'_G}\arrow[r,hook]\&\widehat M_{g,\overline I}
\end{tikzcd}}
où
\[
\text{(59)}\qquad
\left\{\begin{array}{l}
\Gamma^!\subset\Gamma=\operatorname{Aut}(\underline G)\\
\Vert_{\mathrm{def}}\\
\{\deleted{\unsure{g}}u\in\Gamma\mid u_I=\operatorname{id}_I\}=\operatorname{Ker}(\operatorname{Aut}(\underline G)\to\unsure{\mathfrak S}_I),
\end{array}\right.
\]
où $I$ est l'ens. des ``pts marqués'' de $\underline G$.

Alors les \inserted{$2^\mu$} générisations de $\underline G$ sont en correspondance biunivoque avec l'ens. des sous-multiplicités (fermées, \deleted{lisses} donc propres sur $\mathbf Z$, et lisses sur $\mathbf Z$) de $\widehat M_{g,\nu}$ (resp. de $\widehat M_{g,\overline I}$) de la forme $\widehat M_{[\underline G']}$ (resp. $(\widehat M_{\underline G'},\Gamma'^!)$) qui contiennent $\widehat M_{[\underline G]}$ (resp. $(\widehat M_G,\Gamma^!)$). C'est
\margin{49}
